% Generated draft: requires mathematical and source review.
\documentclass[11pt]{article}
\usepackage[margin=25mm]{geometry}
\usepackage{fontspec}
\setmainfont{Latin Modern Roman}
\usepackage{amsmath,amssymb,mathtools,unicode-math}
\setmathfont{Latin Modern Math}
\usepackage{enumitem,needspace,xurl,hyperref}
\setlength{\parindent}{0pt}
\setlength{\parskip}{5pt}
\setlength{\emergencystretch}{3em}
\newcommand{\N}{\mathbb{N}}
\newcommand{\Z}{\mathbb{Z}}
\newcommand{\Q}{\mathbb{Q}}
\newcommand{\R}{\mathbb{R}}
\newcommand{\C}{\mathbb{C}}
\newcommand{\F}{\mathbb{F}}
\newcommand{\A}{\mathbb{A}}
\newcommand{\PP}{\mathbb{P}}
\newcommand{\E}{\mathbb{E}}
\newcommand{\eps}{\varepsilon}
\newcommand{\dd}{\,\mathrm{d}}
\begin{document}
\section*{UnsolvedMath ID: 30005509}
\textbf{Source problem code:} catalogue entry
\subsection{Definitions and mathematical statement}
Let \(p\) be a prime, let \(G\) be a finite group, and let \(k\) be a field of characteristic \(p\). The supplied statement imposes no further restriction on \(G\) or \(k\). Write \(kG\) for the group algebra and \(\operatorname{StMod}(kG)\) for the stable category of all left \(kG\)-modules: its morphisms are module homomorphisms modulo those factoring through projective modules. Write \(\operatorname{stmod}(kG)\) for the full subcategory represented by finitely generated \(kG\)-modules. A thick subcategory is a full triangulated subcategory closed under direct summands. For a collection \(\mathcal S\) of objects, \(\operatorname{thick}_{\operatorname{StMod}(kG)}(\mathcal S)\) denotes the smallest thick subcategory containing \(\mathcal S\).

\textbf{Qualification concerning cohomology.} The supplied reconstruction does not specify the cohomology theory, its degree range, or the coefficient ring over which finite generation is required. Consequently, the following records the intended assertion as a parameterized formulation; those choices must be recovered from the cited source before it becomes an unambiguous conjecture.

Fix a graded coefficient ring \(R\) and a graded cohomology functor \(\mathsf H\) on \(\operatorname{StMod}(kG)\), with values in graded \(R\)-modules, representing the intended conventions. Define
\[
\mathcal C_{\mathrm{fg}}(\mathsf H,R)
 =\{M\in\operatorname{StMod}(kG):\mathsf H(M)\text{ is finitely generated over }R\},
\qquad
\mathcal C_0(\mathsf H)
 =\{M\in\operatorname{StMod}(kG):\mathsf H(M)=0\}.
\]
Here finite generation means generation by finitely many homogeneous elements, and \(\mathsf H(M)=0\) means vanishing in every degree included in the chosen convention.

\textbf{Requested assertion, subject to these unresolved conventions.} For every finite group \(G\) and every field \(k\) of characteristic \(p\), the intended subcategories satisfy
\[
\mathcal C_{\mathrm{fg}}(\mathsf H,R)
 =\operatorname{thick}_{\operatorname{StMod}(kG)}
 \bigl(\operatorname{stmod}(kG)\cup\mathcal C_0(\mathsf H)\bigr).
\]
Thus every module with finitely generated cohomology should be obtainable from finitely generated modules and modules with vanishing cohomology by finitely many triangulated operations and passage to direct summands. The supplied wording also asserts that the left-hand side and the subcategory of modules with vanishing cohomology are thick; this requires the intended cohomology conventions.

For comparison, ordinary group cohomology is \(H^n(G,M)=\operatorname{Ext}_{kG}^n(k,M)\) for \(n\geq0\), with its natural action by \(H^*(G,k)\). Tate cohomology \(\widehat H^n(G,M)\) is indexed by all integers and is defined using a complete projective resolution of the trivial module \(k\). These conventions cannot be silently interchanged. In particular, ordinary \(H^0(G,M)=M^G\) need not vanish on projective modules, so ordinary cohomology including degree zero does not directly define a functor on the stable category. Positive-degree ordinary cohomology avoids that particular issue, but selecting it, or selecting Tate cohomology and a coefficient ring, would add information absent from the supplied record.
\subsection{Short English statement}
The intended conjecture says that every module with finitely generated cohomology belongs to the thick subcategory generated by finitely generated modules and modules with no cohomology. The supplied record leaves the cohomology theory, degree range, and finite-generation convention unresolved, so the displayed equality is conditional on specifying them. Its open status and claimed source verification are supplied curation labels, not independently verified conclusions.
\subsection{Sources}
\begin{itemize}
\item Local UnsolvedMath registry, \texttt{data/unsolvedmath/problems.json}, supplied record with permanent ID \(30005509\), title \emph{Finitely Generated Cohomology in Stable Module Categories}, exact source code \texttt{catalogue entry}. The record belongs to the Oberwolfach Reports Open Problems collection and gives proposed year 2023; no individual proposer is supplied.
\item \emph{Representations of Finite Groups} (2023), Oberwolfach Reports, DOI 10.4171/owr/2023/19, pp. 1031--1084. Supplied report provenance: \url{https://doi.org/10.4171/owr/2023/19}. The background describes Oberwolfach Reports as documenting mathematical workshops and their open-problem discussions.
\item Benson--Carlson, \emph{Modules with finitely generated cohomology}. Supplied literature reference: \url{https://doi.org/10.1016/j.jalgebra.2024.07.040}.
\item Benson--Greenlees, \emph{modules and singularities of C-star-BG}. Title and attribution as supplied; literature reference: \url{https://arxiv.org/abs/2305.08580}.
\item Supplied literature assessment dated \textbf{2026-08-22}: the record labels the problem open, relates the converse to Benson--Greenlees thick-subcategory conjectures, and reports proofs for the \(p\)-nilpotent-centralizer case and a few additional groups, rather than all finite groups. This is supplied dated curation derived from the Oberwolfach material, not an independently verified literature assessment. The record expressly requests independent verification of time-sensitive status and cited claims.
\item The record labels the reconstructed statement \texttt{corrected\_verified} and says it was checked against the official Oberwolfach report. These are supplied verification claims; no source consultation or mathematical verification is asserted here.
\end{itemize}
\end{document}
